\chapter{Solution Design}
\label{ch:solution-design}

% ─────────────────────────────────────────────────────────────────────────────
\section{Functional and Non-Functional Requirements}
\label{sec:requirements}

\subsection*{Functional Requirements}

\begin{description}

\item[FR-1] The system shall allow users to register, authenticate, and access
a personal dashboard displaying their assessment history and current profile.

\item[FR-2] The system shall administer three distinct assessment types:
aptitude tests (multiple-choice, auto-scored), Big Five personality scales
(Likert sliders), and open-ended reflective questions (free text).

\item[FR-3] For open-text responses, the system shall apply NLP processing ---
including tokenisation, lemmatisation, and BERT-based semantic embedding ---
to extract thematic vectors that contribute to the user's vocational profile.

\item[FR-4] The system shall generate a career match score for each candidate
occupation using an explicit, transparent weighted formula:
%
\begin{equation}
  \text{score}(c,\,u)
  \;=\;
  w_1 \cdot \text{aptitude}(u)
  \;+\; w_2 \cdot \text{personality\_fit}(c,u)
  \;+\; w_3 \cdot \text{nlp\_similarity}(c,u)
  \;+\; w_4 \cdot \text{market\_demand}(c)
  \label{eq:score}
\end{equation}
%
where weights are initialised at $(w_1, w_2, w_3, w_4) = (0.25,\,0.30,\,0.25,\,0.20)$,
must sum to $1.0$, and are constrained to the range $[0.05,\,0.50]$ per weight.

\item[FR-5] Each recommendation shall be accompanied by a SHAP-style factor
decomposition showing each weight's percentage contribution to the final score,
expressed in plain language (e.g.\ \emph{``Recommended primarily due to high
Conscientiousness score and strong market demand in this sector''}).

\item[FR-6] Users shall be able to accept or reject each recommendation.
Accepted recommendations shall trigger an incremental weight update: the weight
of the highest-contributing factor for that recommendation is increased by a
learning rate $\alpha = 0.025$, followed by re-normalisation across all weights.

\item[FR-7] The system shall periodically query the JobDataPool API
(\texttt{GET /v1/jobs}, filtered by industry and country code) to maintain an
up-to-date job listings cache. Job listings shall be displayed in the user
dashboard, matched against the user's top recommended career categories.

\item[FR-8] An administrative bias audit module shall compute and display
fairness metrics --- including disparate impact ratio and equal opportunity
scores --- across four protected attributes: gender, age group, field of study,
and socioeconomic background. These attributes shall be collected at onboarding
but excluded from the scoring formula.

\end{description}

\subsection*{Non-Functional Requirements}

\begin{description}

\item[NFR-1] The API shall respond within 800\,ms for standard operations;
NLP inference for open-text responses may take up to 1\,500\,ms and shall be
handled asynchronously with a progress indicator.

\item[NFR-2] All user data shall be encrypted at rest (AES-256) and in transit
(TLS~1.3). Passwords shall be hashed with bcrypt (cost factor $\geq 12$).

\item[NFR-3] The recommendation formula, weight values, and factor contributions
shall be visible to the user at all times. No output shall be presented without
an accompanying explanation.

\item[NFR-4] The interface shall comply with WCAG~2.1~AA accessibility
standards and be responsive across desktop and mobile devices.

\item[NFR-5] The system shall support a minimum of 10\,000 concurrent users
through stateless horizontal scaling of the FastAPI backend.

\item[NFR-6] The bias audit module shall flag any demographic group whose
disparate impact ratio falls below 0.80 (the standard \emph{four-fifths rule})
and generate a written alert for administrative review.

\end{description}

% ─────────────────────────────────────────────────────────────────────────────
\section{Design Objectives}
\label{sec:design-objectives}

The platform is designed around six primary objectives, each traceable to the
requirements established in Section~\ref{sec:requirements} and each reflecting
a deliberate response to the limitations identified in existing career guidance
systems.

\paragraph{DO-1 --- Transparent, auditable scoring.}
The recommendation engine shall expose its internal logic completely. The
weighted formula~\eqref{eq:score} is not an implementation detail but a
user-facing feature: weights are displayed, their current values are shown, and
the contribution of each factor to every recommendation is broken down
explicitly (FR-5). This objective directly addresses the black-box criticism
levelled at systems such as Jobiri and C3-IoC, where users receive
recommendations with no insight into why those recommendations were generated.
Transparency is a design constraint, not an optional layer added afterward.

\paragraph{DO-2 --- Selective application of NLP.}
Natural language processing is computationally expensive and introduces
interpretability challenges. It shall therefore be applied only where linguistic
ambiguity genuinely exists. Aptitude MCQ scoring and Big Five personality
scaling are deterministic operations that require no NLP. NLP is reserved for
two components: semantic analysis of open-text reflective responses (using BERT
embeddings and thematic classification trained on the Q\&A dataset), and skill
normalisation in job matching (mapping free-form skill terms from job listings
to canonical skill labels in the user profile).

\paragraph{DO-3 --- A learning system, not a static one.}
The recommendation model shall improve over time through an explicit feedback
loop (FR-6). When a user accepts a recommendation, the weight of the
highest-contributing scoring factor is nudged upward by $\alpha = 0.025$ and
the weight vector is re-normalised. Over many users, this mechanism allows the
weight vector to converge toward empirically validated distributions rather than
relying on manually initialised defaults. The feedback mechanism is visible:
users can observe weight values change in response to their choices.

\paragraph{DO-4 --- Fairness as a first-class property.}
Four protected demographic attributes are collected at onboarding: gender, age
group, field of study, and socioeconomic background. These are architecturally
isolated --- stored in a separate collection, never passed to the scoring
pipeline, and used exclusively by the bias audit module (FR-8). The audit
module computes disparate impact ratios and equal opportunity scores across all
four axes, flagging violations of the four-fifths rule for administrative review
(NFR-6). Fairness is designed into the data model from the beginning.

\paragraph{DO-5 --- Live labour market grounding.}
Career recommendations shall be connected to real, current job market data
rather than static occupation databases. The JobDataPool API is queried
periodically in a background sync job that populates a local MongoDB cache
of job listings filtered by industry and country. This design separates the
latency-sensitive recommendation pipeline from the network-dependent data
collection process, ensuring the 800\,ms response target (NFR-1) is met while
job data remains reasonably fresh.

\paragraph{DO-6 --- Scalability and maintainability through modularity.}
Each functional component --- assessment engine, NLP processor, recommendation
engine, job aggregator, bias auditor --- is implemented as an independently
deployable service in its own Docker container. This modular architecture
allows individual components to be updated, replaced, or scaled horizontally
without affecting the rest of the system (NFR-5). The NLP inference service,
which has higher computational requirements, can be scaled independently of
the lighter API routes.

% ─────────────────────────────────────────────────────────────────────────────
\section{General Architecture of the Platform}
\label{sec:architecture}

The platform follows a modular, service-oriented architecture in which each
functional concern is encapsulated in an independently deployable component.
Components communicate over HTTP (REST) internally, are containerised via
Docker, and are orchestrated through Kubernetes. This design reflects DO-6:
any single component can be scaled, updated, or replaced without cascading
changes to the rest of the system. The architecture is organised into six layers.

\subsection*{Layer 1 --- Frontend (Vue.js + Tailwind CSS)}

The client application handles all user-facing interactions: authentication,
assessment delivery, results display, recommendation browsing, and the bias
audit dashboard. Vue.js was selected over Angular for its lower configuration
overhead and over React for its single-file component model, which keeps
assessment logic, styling, and template co-located. Tailwind CSS handles
responsive layout without a custom design system. The frontend communicates
exclusively with the Backend API layer.

Two views deserve specific mention. The \emph{Recommendation Dashboard} renders
both the ranked career list and the live weight panel (DO-1), pulling the
current weight vector from the backend on every load so that feedback-driven
changes are always reflected. The \emph{Bias Audit View} is a read-only
reporting interface that queries a dedicated audit endpoint, keeping audit data
flow separate from the recommendation data flow.

\subsection*{Layer 2 --- Backend API (FastAPI, Python)}

The central orchestration layer. FastAPI was chosen for its native async
support, automatic OpenAPI/Swagger documentation generation, and performance
characteristics comparable to Node.js. It exposes the following route groups:

\begin{itemize}
  \item \texttt{/auth} --- registration, login, JWT issuance and refresh
  \item \texttt{/survey} --- assessment session management, answer submission,
        auto-scoring for MCQ types
  \item \texttt{/profile} --- retrieval of the user's aggregated psychometric
        profile and current weight vector
  \item \texttt{/recommendations} --- career match score computation, SHAP
        breakdown generation, recommendation retrieval
  \item \texttt{/feedback} --- acceptance and rejection event ingestion, weight
        update triggering
  \item \texttt{/jobs} --- querying the local job listings cache with skill and
        location filters
  \item \texttt{/audit} --- bias metric computation and report retrieval
        (admin-only, JWT role check)
\end{itemize}

The backend does not run NLP inference directly. Open-text responses are
forwarded asynchronously to the NLP Service (Layer~3) and results are stored
in MongoDB before the recommendation pipeline reads them.

\subsection*{Layer 3 --- NLP Service (HuggingFace Transformers + scikit-learn + NLTK)}

A dedicated Python microservice responsible for all language processing tasks.
It is the only component with a GPU dependency and is therefore scaled
independently. It exposes two internal endpoints:

\begin{itemize}
  \item \texttt{/embed} --- accepts raw open-text input, returns a BERT
        embedding vector using \texttt{sentence-transformers/all-MiniLM-L6-v2}
  \item \texttt{/classify} --- applies a trained SVM classifier over the
        embedding to assign thematic labels used in career matching
\end{itemize}

NLTK handles preprocessing before the transformer model. Scikit-learn hosts
the lightweight classifiers operating on the embeddings. The NLP Service is
called only when open-text responses are present or when skill normalisation
is required for job listing data.

\subsection*{Layer 4 --- Recommendation and Explanation Engine}

Logically part of the backend but described separately because it encompasses
the three core improvements of this work. It consists of three tightly coupled
components: the \emph{Scoring Module} (implements formula~\eqref{eq:score}),
the \emph{SHAP Explanation Module} (computes per-factor contribution
percentages), and the \emph{Feedback and Weight Update Module} (handles
incremental learning). These are described in full in
Section~\ref{sec:modules}.

\subsection*{Layer 5 --- Job Data Sync (JobDataPool API + Python scheduler)}

A background service that runs on a configurable schedule (default: every six
hours) and queries the JobDataPool API (\texttt{GET /v1/jobs}, base URL
\texttt{https://api.jobdatapool.com}) filtered by industry and country code.
Retrieved listings are normalised via the NLP Service skill normalisation
pipeline and stored in a MongoDB \texttt{jobs\_snapshot} collection. The rate
limit of 30 requests per 60-second window is respected by introducing a
2-second delay between industry queries within a single sync cycle.

\subsection*{Layer 6 --- Data Layer (MongoDB)}

MongoDB stores all persistent state: user profiles, assessment results, NLP
analysis outputs, weight vectors, recommendation history, job listings cache,
feedback events, session tokens for refresh-token revocation, and bias audit
snapshots. The document-oriented model accommodates the heterogeneous structure
of assessment responses without requiring schema migrations as new test types
are introduced.

\begin{figure}[h]
  \centering
  % Insert architecture diagram here (generated from Section 3.3 design)
  \fbox{\parbox{0.85\textwidth}{\centering\vspace{2cm}
    Figure 3.1: General architecture diagram\\
    (Vue.js Frontend $\rightarrow$ FastAPI Backend $\rightarrow$
    NLP Service / Recommendation Engine / Job Sync $\rightarrow$ MongoDB)
  \vspace{2cm}}}
  \caption{General architecture of the platform.}
  \label{fig:architecture}
\end{figure}

% ─────────────────────────────────────────────────────────────────────────────
\section{Functional Modules}
\label{sec:modules}

\subsection*{Module 1 --- User Management}

Users register with a valid email address and a password hashed using bcrypt
(cost factor~12). On login, the backend issues a short-lived JWT access token
(15-minute expiry) and a long-lived refresh token (7-day expiry) stored
server-side in a \texttt{sessions} collection, enabling session continuity
with the ability to revoke sessions explicitly.

The user document in MongoDB stores two logically separate categories of data.
The first is the \emph{psychometric profile}: aptitude scores, Big Five
personality vector, NLP thematic labels, and the current weight vector. The
second is the \emph{demographic profile}: gender, age group, field of study,
and socioeconomic background. These two categories are stored in separate
MongoDB collections --- \texttt{users} and \texttt{user\_demographics} --- and
joined only within the bias audit pipeline. The recommendation engine reads
exclusively from \texttt{users}, enforcing at the data layer the guarantee that
protected attributes never influence scoring.

Users may request full account deletion at any time, triggering cascade removal
from both collections in compliance with GDPR Article~17.

\subsection*{Module 2 --- Assessment Engine}

The assessment engine delivers three test types, each with a distinct scoring
mechanism. The key design principle is that scoring complexity matches the
nature of the input --- no NLP is applied where deterministic scoring suffices.

\paragraph{Aptitude tests (MCQ).}
Two datasets supply the question bank: the ScienceQA dataset (HuggingFace) for
scientific reasoning, geography, and language questions, and the Engineering
Aptitude dataset (Kaggle) for numerical and logical-mathematical reasoning.
Both datasets are exclusively multiple-choice with integer answer indices.
Scoring is deterministic:
%
\begin{equation}
  \text{aptitude\_score}(u)
  \;=\;
  \frac{\text{correct\_answers}}{\text{total\_questions}} \times 100
  \label{eq:aptitude}
\end{equation}
%
No NLP is involved. The result is a single normalised integer in $[0, 100]$
that feeds directly into formula~\eqref{eq:score} as the aptitude factor.

\paragraph{Big Five personality scales.}
Each of the five OCEAN dimensions is presented as a labelled slider. The user's
position is recorded as a float in $[0, 1]$. Scoring is deterministic --- the
slider value is the score. The five values form the personality vector
$[O, C, E, A, N]$, used in career-personality matching via cosine similarity
against pre-computed career archetype vectors derived from the OCEAN dataset.
No NLP is involved.

\paragraph{Open-text reflective questions.}
The only assessment type requiring NLP processing. The user's free-text
response is passed to the NLP Service and returns as a thematic embedding
vector and a set of classified labels stored in the \texttt{nlp\_analysis}
collection. Submission is handled asynchronously: the frontend shows a
progress indicator while the backend awaits the NLP Service response.

\subsection*{Module 3 --- NLP Processing Service}

The NLP Service processes two distinct inputs: open-text assessment responses
and job listing skill fields from the JobDataPool sync. The processing pipeline
consists of three stages.

\paragraph{Stage 1 --- Preprocessing (NLTK).}
Raw text is lowercased, tokenised, stripped of stopwords, and lemmatised.
For open-text responses this reduces noise from grammatical variation. For
job skill strings it additionally handles tokenisation of compound terms.

\paragraph{Stage 2 --- Semantic embedding (HuggingFace Transformers).}
The preprocessed text is passed to
\texttt{sentence-transformers/all-MiniLM-L6-v2}, producing a 384-dimensional
dense vector. This model is chosen over \texttt{bert-base-uncased} for two
reasons: it runs inference approximately four times faster, and its
sentence-level training objective is more appropriate for paragraph-length
inputs. For the prototype scope, MiniLM provides an acceptable
accuracy-latency trade-off.

\paragraph{Stage 3 --- Thematic classification (scikit-learn SVM).}
The 384-dimensional embedding is passed to an SVM classifier trained on the
Q\&A dataset (Kaggle). The classifier assigns one or more thematic labels from
a predefined taxonomy: \texttt{analytical}, \texttt{creative},
\texttt{interpersonal}, \texttt{technical}, \texttt{leadership},
\texttt{structured}. These labels, together with the raw embedding vector, are
stored in \texttt{nlp\_analysis} and used downstream in two ways: the embedding
feeds into career similarity scoring, and the labels feed into SHAP explanation
generation.

\paragraph{Skill normalisation (job matching only).}
When processing job listing skill fields from JobDataPool, the NLP Service
additionally maps free-form skill terms to a canonical skill vocabulary. For
example, \emph{pandas}, \emph{data wrangling}, and \emph{dataframe
manipulation} are all mapped to the canonical label \emph{Data Analysis}. This
enables meaningful matching between user profile labels and job skill
requirements without relying on exact string matching.

\subsection*{Module 4 --- Recommendation and Explanation Engine}

\paragraph{4a --- Weighted Scoring.}
Career match scores are computed using formula~\eqref{eq:score}. The four
factor functions are defined as follows:

\begin{itemize}
  \item $\text{aptitude}(u)$ --- the user's normalised aptitude score in
        $[0, 1]$
  \item $\text{personality\_fit}(c, u)$ --- cosine similarity between the
        user's OCEAN vector and the archetype OCEAN vector for career $c$,
        mapped to $[0, 1]$
  \item $\text{nlp\_similarity}(c, u)$ --- cosine similarity between the
        user's NLP embedding vector and the centroid embedding vector for
        career $c$, computed from training examples in the Q\&A dataset
  \item $\text{market\_demand}(c)$ --- a normalised signal derived from the
        frequency and recency of job listings for career $c$ in the
        JobDataPool cache, scaled to $[0, 1]$
\end{itemize}

Weights are initialised at $(w_1, w_2, w_3, w_4) = (0.25,\,0.30,\,0.25,\,0.20)$,
constrained to $[0.05,\,0.50]$ per weight, and normalised to sum to $1.0$ at
all times. The current weight vector is stored per user in the \texttt{weights}
document and updated by the feedback component.

\paragraph{4b --- SHAP-style Explanation.}
Because the scoring model is a weighted linear sum, the contribution of each
factor to the final score is analytically exact. The percentage contribution
of factor $i$ to career $c$'s score is:
%
\begin{equation}
  \text{contribution}_i(c, u)
  \;=\;
  \frac{w_i \cdot f_i(c, u)}{\text{score}(c, u)} \times 100
  \label{eq:shap}
\end{equation}
%
These four values sum to 100\% and are attached to each recommendation as a
structured breakdown. A plain-language summary is generated alongside ---
for example: \emph{``Recommended primarily due to strong personality fit
(41\%) and high market demand in this sector (28\%). Your text responses
showed strong analytical and technical orientation, contributing 22\% to
this match.''}

\paragraph{4c --- Incremental Weight Update (Feedback Loop).}
When a user accepts recommendation $c^*$, the following update is applied:
%
\begin{align}
  i^* &= \arg\max_i\; \text{contribution}_i(c^*, u)
  \label{eq:argmax}\\
  w_{i^*} &\leftarrow w_{i^*} + \alpha
  \label{eq:update}\\
  \mathbf{w} &\leftarrow \frac{\mathbf{w}}{\sum_j w_j}
  \qquad \text{(re-normalisation)}
  \label{eq:renorm}
\end{align}
%
where $\alpha = 0.025$. When a user rejects a recommendation, no weight update
is applied; the rejection event is logged in the \texttt{feedback} collection
for aggregate analysis. The top five careers by score are returned, each with
its SHAP breakdown, plain-language explanation, salary range, and matched live
job listings.

\subsection*{Module 5 --- Job Data Aggregation}

This module replaces web scraping with a structured API integration. The
JobDataPool API is queried by a Python background scheduler running as a
separate containerised service. Queries are parameterised by industry and
country code. Each sync cycle executes the following steps:

\begin{enumerate}
  \item For each industry in the platform's career taxonomy, issue a query to
        \texttt{/v1/jobs?industries=\{industry\}\&country\_code=\{code\}\&limit=250}
  \item Pass returned skill strings through the NLP Service skill normalisation
        pipeline
  \item Upsert normalised records into the \texttt{jobs\_snapshot} MongoDB
        collection, keyed on \texttt{job\_id}
  \item Record a sync timestamp; listings older than 14 days are marked stale
        and excluded from dashboard queries
\end{enumerate}

The sync cycle runs every six hours. The rate limit of 30 requests per
60-second window is respected by introducing a 2-second delay between industry
queries within a single cycle.

% ─────────────────────────────────────────────────────────────────────────────
\section{Data Flow}
\label{sec:dataflow}

The platform operates across three concurrent data flows. They share the same
MongoDB instance but are otherwise independent.

\subsection*{Flow 1 --- Assessment to Recommendation}

\begin{enumerate}
  \item \textbf{Input collection.} The user completes the three assessment
        types. MCQ answers and personality slider values are submitted to the
        backend via \texttt{/survey/submit}. The open-text response is submitted
        separately and flagged for NLP processing.

  \item \textbf{Deterministic scoring.} The backend scores the aptitude MCQ
        and records Big Five slider values directly. These results are persisted
        to the \texttt{results} collection synchronously.

  \item \textbf{Asynchronous NLP processing.} The open-text response is
        forwarded to the NLP Service. The backend returns HTTP~202 to the
        frontend, which displays a progress indicator. The NLP Service runs the
        three-stage pipeline and writes the embedding vector and thematic labels
        to \texttt{nlp\_analysis}.

  \item \textbf{Profile assembly.} Once all inputs are scored, the backend
        assembles the full user profile and persists it to \texttt{users}.

  \item \textbf{Scoring and explanation.} The recommendation engine applies
        formula~\eqref{eq:score} across all candidate careers, ranks the
        results, computes SHAP breakdowns via equation~\eqref{eq:shap}, and
        attaches matching job listings from \texttt{jobs\_snapshot}. The top
        five results are persisted to \texttt{recommendations}.

  \item \textbf{Delivery.} The frontend retrieves results via
        \texttt{/recommendations} and renders the career list, SHAP
        breakdowns, and weight panel.
\end{enumerate}

\subsection*{Flow 2 --- Feedback and Weight Update}

\begin{enumerate}
  \item \textbf{Event capture.} The user accepts or rejects a recommendation.
        The frontend sends a POST to \texttt{/feedback}.

  \item \textbf{Weight update (accept only).} The backend retrieves the stored
        SHAP breakdown, identifies the highest-contributing factor, and applies
        equations~\eqref{eq:update}--\eqref{eq:renorm}. The updated weight
        vector is persisted to the \texttt{weights} document.

  \item \textbf{Rejection logging.} If the action is a rejection, the event is
        written to \texttt{feedback} with a timestamp and factor breakdown.

  \item \textbf{Live weight refresh.} The updated weight vector is returned in
        the API response. The frontend updates the weight panel immediately.
\end{enumerate}

\subsection*{Flow 3 --- Background Job Sync}

\begin{enumerate}
  \item \textbf{Scheduled trigger.} The Python scheduler fires every six hours
        and iterates over the platform's career taxonomy.

  \item \textbf{Skill normalisation.} Retrieved listings are passed through
        the NLP Service skill normalisation pipeline.

  \item \textbf{Cache upsert.} Normalised listings are upserted into
        \texttt{jobs\_snapshot}. Listings older than 14 days are marked stale.

  \item \textbf{Bias snapshot (on demand).} When an administrator triggers a
        bias report via \texttt{/audit}, a read-only pipeline joins
        \texttt{recommendations} with \texttt{user\_demographics}, computes
        fairness metrics, and writes a snapshot to \texttt{bias\_audit}. This
        join is the only point in the system where demographic data and
        recommendation data are read together.
\end{enumerate}

\begin{figure}[h]
  \centering
  % Insert data flow diagram here (generated from Section 3.5 design)
  \fbox{\parbox{0.85\textwidth}{\centering\vspace{2cm}
    Figure 3.2: Three concurrent data flows\\
    (Flow 1: Assessment $\to$ Recommendation;\;
     Flow 2: Feedback $\to$ Weight update;\;
     Flow 3: Job sync $\to$ Cache)\\
    Shared MongoDB at convergence point.
  \vspace{2cm}}}
  \caption{Concurrent data flows in the platform.}
  \label{fig:dataflow}
\end{figure}

% ─────────────────────────────────────────────────────────────────────────────
\section{Datasets Used}
\label{sec:datasets}

The platform draws on five open static datasets and one live API source. Each
was selected to cover a specific functional component.

\subsection{Dataset Overview and Integration Rationale}

\begin{table}[h]
\centering
\caption{Data sources, roles, and processing modes.}
\label{tab:datasets}
\renewcommand{\arraystretch}{1.35}
\begin{tabular}{p{3.4cm} p{1.5cm} p{2.6cm} p{1.5cm} p{3.6cm}}
\hline
\textbf{Source} & \textbf{Type} & \textbf{Module} & \textbf{NLP} & \textbf{Primary role} \\
\hline
ScienceQA (HuggingFace)
  & Static & Assessment engine & No
  & MCQ aptitude items --- science, geography, language \\
Engineering Aptitude (Kaggle)
  & Static & Assessment engine & No
  & MCQ aptitude items --- numerical and logical reasoning \\
Q\&A Dataset (Kaggle)
  & Static & NLP service & Yes
  & Trains the SVM thematic classifier \\
OCEAN Big Five (HuggingFace)
  & Static & Recommendation engine & No
  & Derives career archetype personality vectors \\
AI-Powered Job Market (Kaggle)
  & Static & Recommendation engine & Partial
  & Initialises market demand signals and career taxonomy \\
JobDataPool API
  & Live API & Job sync service & Partial
  & Live job listings cache, refreshed every 6 hours \\
\hline
\end{tabular}
\end{table}

Three points in Table~\ref{tab:datasets} warrant explanation.

First, the OCEAN dataset is marked as not requiring NLP. OCEAN scores are
numerical self-reported values --- no language processing is needed. They are
used to compute career archetype vectors: for each career category, the mean
OCEAN vector across users in that occupation becomes the reference vector
against which user profiles are compared via cosine similarity.

Second, the AI-Powered Job Market Kaggle dataset and the JobDataPool API serve
complementary rather than redundant roles. The Kaggle dataset is static and
used at initialisation: it provides the career taxonomy, seeds the market
demand signals used in formula~\eqref{eq:score}, and supplies training
examples for the career matching component. JobDataPool is used at runtime:
it provides live job listings and fresh market demand signals.

Third, skill normalisation is listed as \emph{partial} NLP for both
job-related sources. This refers specifically to the skill normalisation
pipeline in the NLP service --- lightweight similarity matching, not full
semantic embedding.

\subsection{Per-Dataset Integration Detail}

\paragraph{ScienceQA and Engineering Aptitude (MCQ datasets).}
Both datasets supply the question bank for the aptitude assessment module.
ScienceQA~\cite{scienceqa} covers scientific reasoning, geography, and
language questions across grades 2--11 with 21\,264 items in total. All items
are multiple-choice with a predetermined \texttt{choices} array and an integer
\texttt{answer} index. Scoring is purely deterministic; no NLP is involved.
The \texttt{lecture} and \texttt{solution} fields, unique to ScienceQA, are
used to provide post-answer explanations, supporting the explainability
objective (DO-1). Engineering Aptitude provides numerical sequences, spatial
reasoning, and quantitative problem-solving items, complementing ScienceQA's
science and language coverage.

\paragraph{Q\&A Dataset (NLP classifier training).}
Used offline to train the SVM thematic classifier. Items are question-answer
pairs with sufficient textual content for reliable thematic labelling. The
classifier is trained to assign labels from the vocational taxonomy
(\texttt{analytical}, \texttt{creative}, \texttt{interpersonal},
\texttt{technical}, \texttt{leadership}, \texttt{structured}) based on BERT
embeddings. The trained classifier is serialised and loaded by the NLP service
at startup.

\paragraph{OCEAN Big Five Dataset.}
Used in two ways. First, it calibrates the normalisation of slider values so
that a score of $0.7$ on Conscientiousness corresponds to a meaningful
percentile. Second, for occupational categories with sufficient representation
($n \geq 30$), the mean OCEAN vector per occupation becomes the career
archetype vector stored in \texttt{career\_archetypes}. For categories with
$n < 30$, archetype vectors are initialised from published RIASEC-to-Big-Five
mappings as a validated theoretical fallback.

\paragraph{AI-Powered Job Market Dataset (Kaggle).}
Provides the initial career taxonomy, seeds the market demand signals, and
supplies training examples for career matching. Skill strings are passed
through the NLP Service skill normalisation pipeline. Processed records seed
the \texttt{career\_archetypes} and \texttt{jobs\_snapshot} collections at
initialisation.

\paragraph{JobDataPool API (live source).}
Queried on a background schedule as described in Section~\ref{sec:modules}
(Module~5). Populates and maintains \texttt{jobs\_snapshot} with current job
listings and provides fresh market demand signals that periodically recalibrate
the static initialisation values.

\subsection{Preprocessing and Filtering}

All static datasets undergo a common preprocessing pipeline before being loaded
into the platform.

\paragraph{MCQ datasets (ScienceQA and Engineering Aptitude).}
Duplicate questions are removed by exact text match. For ScienceQA, items
where \texttt{image} is not null are excluded (text-only subset), and items
from grades below 7 are filtered out to match the university student target
audience. The resulting question pool is serialised as a structured JSON array
and stored in the \texttt{surveys} MongoDB collection at initialisation.
ScienceQA's \texttt{lecture} and \texttt{solution} fields are preserved in the
stored objects and returned to the frontend after the user submits an answer.

\paragraph{Q\&A dataset.}
Items shorter than 30 tokens after preprocessing are removed. The remaining
items are lowercased, tokenised, and lemmatised via NLTK. Class balance across
the six thematic labels is checked; overrepresented classes are downsampled to
a maximum of $2\times$ the size of the smallest class. The balanced corpus is
split 80/20 for training and validation.

\paragraph{OCEAN dataset.}
Responses with more than 20\% missing values are dropped. Remaining scores are
normalised to $[0, 1]$ per dimension using min-max scaling. Occupational
labels are standardised to a controlled vocabulary aligned with the career
taxonomy. Archetype vectors are computed as per-occupation means with a minimum
sample size of $n = 30$.

\paragraph{AI-Powered Job Market dataset.}
Duplicate listings are removed. Salary fields are parsed and normalised to
annual EUR ranges; entries with unparseable salary fields are retained but
flagged as salary-unknown. Skill strings are passed through the NLP Service
skill normalisation pipeline. Processed records seed \texttt{career\_archetypes}
and \texttt{jobs\_snapshot} at initialisation.

\paragraph{JobDataPool API responses.}
Handled at runtime by the job sync service. Skill normalisation,
deduplication by \texttt{job\_id}, and stale-record pruning are applied on
each sync cycle.

% ─────────────────────────────────────────────────────────────────────────────
\section{Technical Considerations}
\label{sec:technical}

\subsection{Performance}

The primary performance target is NFR-1: API responses under 800\,ms for
standard operations, with NLP inference allowed up to 1\,500\,ms handled
asynchronously.

\paragraph{MongoDB indexing strategy.}
Compound indexes are defined at initialisation on the fields most commonly used
in query filters: \texttt{user\_id} on \texttt{results}, \texttt{nlp\_analysis},
\texttt{recommendations}, \texttt{weights}, and \texttt{feedback} collections;
\texttt{industry} and \texttt{country\_code} on \texttt{jobs\_snapshot}; and
\texttt{career\_id} on \texttt{career\_archetypes}.

\paragraph{Asynchronous NLP handling.}
The FastAPI backend issues HTTP~202 immediately upon receiving a text submission
and delegates processing to the NLP Service. The frontend polls a status
endpoint at 1-second intervals. This prevents BERT inference latency ---
typically 400--900\,ms on CPU, under 200\,ms on GPU --- from blocking the
main API thread.

\paragraph{NLP result persistence.}
Once computed, a user's NLP embedding and thematic labels are stored in
\texttt{nlp\_analysis} and reused for all subsequent recommendation queries
without re-inference. Re-embedding is only triggered if the user submits a
new open-text response.

\paragraph{Job sync decoupling.}
The JobDataPool sync job runs as a separate containerised process. All
user-facing job queries read from the local \texttt{jobs\_snapshot} cache,
so external API latency has no effect on user-facing response times.

\paragraph{BERT model selection.}
At 22M parameters, \texttt{all-MiniLM-L6-v2} runs approximately four times
faster than \texttt{bert-base-uncased} (110M parameters). On CPU inference
this translates to roughly 80--120\,ms per embedding versus 400--600\,ms,
making the 1\,500\,ms async budget achievable without GPU hardware in the
prototype deployment.

\subsection{Security}

\paragraph{Authentication.}
Passwords are hashed with bcrypt (cost factor~12). Authentication issues a
short-lived JWT access token (15-minute expiry) and a long-lived refresh token
(7-day expiry) stored server-side in \texttt{sessions} for explicit
revocation. All tokens are transmitted over TLS~1.3.

\paragraph{Data encryption at rest.}
Sensitive fields --- hashed passwords, session tokens, and demographic data in
\texttt{user\_demographics} --- are encrypted at the application layer using
AES-256 before being written to MongoDB.

\paragraph{Demographic data isolation.}
The \texttt{user\_demographics} collection is accessible only through the
\texttt{/audit} route group, which requires a verified admin role claim in the
JWT. No other route joins or reads from this collection.

\paragraph{Input validation.}
All incoming request bodies are validated against Pydantic schemas at the
FastAPI layer. String fields are length-bounded and sanitised, preventing
injection attacks from reaching MongoDB or the NLP Service.

\paragraph{Audit logging.}
Every request to the \texttt{/auth}, \texttt{/recommendations},
\texttt{/feedback}, and \texttt{/audit} route groups writes a log entry
containing a timestamp, anonymised IP address (last octet zeroed), operation
type, and HTTP status code. Logs are retained for 30 days.

\subsection{Scalability and DevOps}

\paragraph{Containerisation.}
Each service runs in its own Docker container built from a minimal base image
(\texttt{python:3.11-slim}). Services communicate over an internal Docker
network; only the FastAPI backend exposes an external port.

\paragraph{Horizontal scaling.}
The FastAPI backend is stateless by design: all session state lives in MongoDB
and JWT tokens carry no server-side state beyond the revocation list.
Additional backend instances can be added behind a load balancer without
session affinity requirements. The NLP Service can similarly be scaled
horizontally.

\paragraph{CI/CD pipeline.}
A GitHub Actions workflow runs on every push to the main branch across four
stages: syntax linting (flake8, ESLint), unit tests (pytest, Vitest), Docker
image build, and deployment to the staging environment. A manual approval gate
separates staging from production. Swagger documentation is auto-generated and
published as a build artefact on every successful run.

\paragraph{Database backup.}
MongoDB performs incremental backups daily, retained for 14 days. A full
restore is tested monthly and documented to complete within 30 minutes. Backup
files are encrypted using the same AES-256 key used at the application layer.

\subsection{Known Technical Constraints}

\paragraph{NLP classifier accuracy is bounded by training data.}
The SVM thematic classifier is trained on the Q\&A Kaggle dataset, which
covers general question-answering content and has not been trained on
career-specific language. Users whose open-text responses use domain-specific
vocabulary outside the training distribution may receive less accurate thematic
classifications. Retraining on accumulated platform-specific open-text
responses is the intended mitigation, listed as a future development direction
in Section~\ref{ch:conclusions}.

\paragraph{Career archetype vectors are partially theory-derived.}
For career categories with fewer than 30 OCEAN-labelled examples, archetype
vectors fall back to published RIASEC-to-Big-Five mappings. These mappings are
well-validated in the literature but introduce a degree of theoretical
approximation where empirical data is sparse.

\paragraph{Weight initialisation is not empirically calibrated.}
The default weight vector $(0.25,\,0.30,\,0.25,\,0.20)$ is initialised based
on the relative theoretical importance of each factor in the career counselling
literature rather than from optimisation on labelled outcome data. The feedback
loop is designed to correct this over time, but in early deployment --- before
sufficient feedback has accumulated --- recommendation quality depends on these
initial values being reasonable rather than optimal.

\paragraph{JobDataPool coverage is not exhaustive.}
For niche or highly specialised career categories, listing volumes may be
insufficient to produce reliable market demand signals. In these cases the
market demand factor defaults to the static initialisation value from the
Kaggle job market dataset.

\paragraph{Algorithmic bias has not been formally evaluated in this prototype.}
The bias audit module computes and reports fairness metrics, but no corrective
intervention has been implemented. Detected disparities are flagged for
administrative review but do not trigger automatic model adjustment.
Implementing bias mitigation interventions is identified as a priority future
development direction in Section~\ref{ch:conclusions}.
